% !TEX TS-program = pdflatex
% !TeX program = pdflatex
% !TEX encoding = UTF-8
% !TEX spellcheck = en_US

\documentclass[xcolor=table]{beamer}

\usepackage{../extra/beamer/karimnlp}

\input{options}

\subtitle[03- Basic text processing]{Chapter 03\\Basic text processing} 

\changegraphpath{../img/basic/}

\begin{document}
	
\begin{frame}
	\frametitle{\inserttitle}
	\framesubtitle{\insertshortsubtitle: Introduction}
	
	\begin{itemize}
		\item Text is structured into paragraphs, sentences, words, and characters.
		\item Various NLP tasks utilize this hierarchical structure for analysis.
		\item Words are commonly used as the primary unit of processing.
		\item Preprocessing phase includes:
		\begin{itemize}
			\item Segmenting text into sentences and words;
			\item Removing irrelevant or redundant information (e.g., stopwords, punctuation);
			\item Normalizing words to reduce variation (e.g., lowercasing, stemming, lemmatization).
		\end{itemize}
	\end{itemize}

\end{frame}


\begin{frame}
	\frametitle{\inserttitle}
	\framesubtitle{\insertshortsubtitle: Plan}
	
	\begin{multicols}{2}
	%	\small
	\tableofcontents
	\end{multicols}
\end{frame}

%===================================================================================
\section{Characters}
%===================================================================================

\begin{frame}
	\frametitle{\insertshortsubtitle}
	\framesubtitle{\insertsection}

	\begin{itemize}
		\item \optword{Regular expressions}
		\begin{itemize}
			\item Used to search for patterns or substrings in text.
			\item Can recognize Type-3 languages (regular languages) in the Chomsky hierarchy.
			\item Useful for lexical analysis and word tokenization.
		\end{itemize}
		\item \optword{Edit distance}
		\begin{itemize}
			\item Used to quantify the dissimilarity between two strings.
			\item Useful for approximate string matching (fuzzy search).
		\end{itemize}
	\end{itemize}

\end{frame}

\subsection{Regular expressions}

\begin{frame}
	\frametitle{\insertshortsubtitle: \insertsection}
	\framesubtitle{\insertsubsection: Syntax (1)}

	\SetTblrInner{rowsep=0pt,colsep=1pt}
	\begin{tblr}{
			colspec = {p{.1\textwidth}p{.41\textwidth}p{.45\textwidth}},
			row{odd} = {lightblue, font=\small},
			row{even} = {lightyellow, font=\small},
			row{1} = {darkblue, font=\bfseries},
		} 
		\textcolor{white}{RE} & \textcolor{white}{Meaning} & \textcolor{white}{Example} \\
		
		. & any character & \keyword{beg.n}: I \expword{begun} at the \expword{begin}ning. \\
		
		 [aeuio] & specific characters & \keyword{[Cc][au]t}: \expword{Cat}s are \expword{cut}e. \\
		 
		[a-e] & character range & \keyword{[A-Z]..}: \expword{I s}aw \expword{Kar}im. \\
		
		[\textasciicircum aeuio] & exclude characters & \keyword{[\textasciicircum A-Z]a.}: I \expword{saw} Karim. \\
		
		c? & one or zero & \keyword{colou?r}: It is \expword{colour} or \expword{color}. \\
		
		c* & zero or more & \keyword{Ye*s} : \expword{Ys}! \expword{Yes}! \expword{Yeeees}! \\
		
		c+ & one or more & \keyword{Ye+s} : Ys! \expword{Yes}! \expword{Yeeees}! \\
		
		c\{n\} & n occurrences & \keyword{Ye\{3\}s} : Ys! Yes! Yees! \expword{Yeees}! \\
		
		c\{n,m\} & between n and m occurrences & \keyword{Ye\{1,2\}s} : Ys! \expword{Yes}! \expword{Yees}! Yeees! \\
		
		c\{n,\} & at least n occurrences & \keyword{Ye\{2,\}s} : Ys! Yes! \expword{Yees}! \expword{Yeees}! \\
		
		c\{,m\} & at most m occurrences & \keyword{Ye\{,2\}s} : \expword{Ys}! \expword{Yes}! \expword{Yees}! Yeees! \\
		
	\end{tblr}

\end{frame}

\begin{frame}
	\frametitle{\insertshortsubtitle: \insertsection}
	\framesubtitle{\insertsubsection: Syntax (2)}

	\begin{minipage}{.6\textwidth}
	\begin{itemize}
		\item Grouping: \keyword{( )}
		\begin{itemize}
			\item E.g., \expword{/([Pp]retty, )+/}: \expword{Pretty, pretty, }pretty good.
		\end{itemize}
		\item Disjunction: \keyword{\textbar}
		\begin{itemize}
			\item E.g., \expword{/continu(e\textbar ation\textbar al(ly)?)/}
		\end{itemize}
		\item Text beginning: \keyword{\textasciicircum}
		\begin{itemize}
			\item E.g., \expword{/\textasciicircum K/}:  \expword{K}ey of Karim.
		\end{itemize}
		\item Text ending: \keyword{\$}
		\begin{itemize}
			\item E.g., \expword{/\textbackslash .[\textasciicircum .]+\$/} :  file.tar\expword{.gz}
		\end{itemize}
		\item To capture a group: \keyword{\$n} or \keyword{\textbackslash n} where \expword{n} is the group's position.
		\begin{itemize}
			\item E.g., \expword{/(.*)(ual\textbar uation)\$/\textbackslash 1ue/}.
		\end{itemize}
	\end{itemize}
	\end{minipage}
	\begin{minipage}{.38\textwidth}
	\begin{tblr}{
			colspec = {p{.2\textwidth}p{.6\textwidth}},
			row{odd} = {lightblue},
			row{even} = {lightyellow},
			row{1} = {darkblue},
		} 
	
		\textcolor{white}{RE} & \textcolor{white}{Equivalence} \\
		
		\textbackslash d & [0-9] \\
		\textbackslash D & [\textasciicircum 0-9] \\
		\textbackslash w & [a-zA-Z0-9\_] \\
		\textbackslash W & [\textasciicircum \textbackslash w] \\
		\textbackslash s & [ \textbackslash r\textbackslash t\textbackslash n\textbackslash f] \\
		\textbackslash S & [\textasciicircum \textbackslash s] \\
	\end{tblr}
	
	\end{minipage}

\end{frame}

\begin{frame}
	\frametitle{\insertshortsubtitle: \insertsection}
	\framesubtitle{\insertsubsection: Use case}
	
	\begin{itemize}
		\item Text editors use regular expressions (REs) for search and replacement.
		\item Most programming languages provide built-in support for REs.
		\item Data extraction: e.g., extracting emails and phone numbers from blogs and social media.
		\item \url{https://github.com/kariminf/aruudy}
		\begin{itemize}
			\item Personal project: Arabic poems meter detection.
			\item Pros: rules are readable and interpretable.
			\item Cons: requires multiple passes to process a single verse.
		\end{itemize}
	\end{itemize}

\end{frame}

\begin{frame}
	\frametitle{\insertshortsubtitle: \insertsection}
	\framesubtitle{\insertsubsection: Some humor}

	\begin{center}
		\vgraphpage{humor/humor-regex.jpg}
	\end{center}

\end{frame}

\subsection{Edit distance}

\begin{frame}
	\frametitle{\insertshortsubtitle: \insertsection}
	\framesubtitle{\insertsubsection: Edit operations}

	\begin{itemize}
		\item \optword{Insertion}: adding a character to a string.\\
		$uv \rightarrow uxv \,/\, u, v \in X^*;\, uv \in X^+;\, x \in X$
		\begin{itemize}
			\item \expword{case $ \rightarrow $ casse, entraînement $ \rightarrow $ entraînnement }
		\end{itemize}
		
		\item \optword{Deletion}: removing a character from a string.\\
		$uxv \rightarrow uv \,/\, u, v \in X^*;\, uv \in X^+;\, x \in X$
		\begin{itemize}
			\item \expword{future $ \rightarrow $ futur, héros $ \rightarrow $ héro}
		\end{itemize}
		
		\item \optword{Substitution}: replacing one character with another.\\
		$uxv \rightarrow uyv \,/\, u, v \in X^*;\, x, y \in X;\, x \ne y$
		\begin{itemize}
			\item \expword{design $ \rightarrow $ dezign, croient $ \rightarrow $ croyent }
		\end{itemize}
		
		\item \optword{Transposition}: swapping two characters.\\
		$uxwyv \rightarrow uywxv \,/\, u, v, w \in X^*;\, x, y \in X;\, x \ne y$
		\begin{itemize}
			\item \expword{play $ \rightarrow $ paly, cueillir $ \rightarrow $ ceuillir}
		\end{itemize}
	\end{itemize}

\end{frame}

\begin{frame}
	\frametitle{\insertshortsubtitle: \insertsection}
	\framesubtitle{\insertsubsection: Algorithms}
	
	\begin{itemize}
		\item \optword{Hamming distance}: considers only substitutions. Strings must be of equal length.
		\item \optword{Longest common subsequence (LCS)}: measures similarity by allowing only insertions and deletions.
		\item \optword{Levenshtein distance}: allows insertion, deletion, and substitution.
		\item \optword{Jaro distance}: measures string similarity and is sensitive to transpositions.
		\item \optword{Damerau–Levenshtein distance}: allows insertion, deletion, substitution, and transposition of adjacent characters.
	\end{itemize}

\end{frame}


\begin{frame}
	\frametitle{\insertshortsubtitle: \insertsection}
	\framesubtitle{\insertsubsection: Hamming distance}
	
	\begin{itemize}
		\item $X$: source string of length $n$.
		\item $Y$: target string of length $n$.
	\end{itemize}
	
	\[
	D(X,Y) = \sum_{i=1}^{n} [X_i \ne Y_i]
	\]
	
%	\[
%	D_{i} = D_{i-1} + 
%	\begin{cases}
%	0 \text{ if } X_i = Y_i\\
%	1 \text{ otherwise}
%	\end{cases}
%	\]
	
	\begin{exampleblock}{Examples of Hamming distance}
		\begin{itemize}
			\item D(100100, 101001) = 3
			\item D(abdelkrim, abderahim) = 3
		\end{itemize}
	\end{exampleblock}
	
\end{frame}

\begin{frame}
	\frametitle{\insertshortsubtitle: \insertsection}
	\framesubtitle{\insertsubsection: Levenshtein distance}
	
	\begin{itemize}
		\item $X$: source string of length $n$.
		\item $Y$: target string of length $m$.
		\item $D$: a matrix, where $D[i, j]$ represents the edit distance between prefixes $X[1..i]$ and $Y[1..j]$.
		\item To compute $D[n, m]$, dynamic programming is used.
		\item Base cases:
		\[
		D[i,0] = i, \quad D[0,j] = j, \quad D[0,0] = 0
		\]
	\end{itemize}
	
	\[
	D[i, j] = \min 
	\begin{cases}
		D[i - 1, j] + 1 & \text{(Deletion)}\\
		D[i, j-1] + 1 & \text{(Insertion)}\\
		D[i-1, j-1] + 
		\begin{cases}
			1 & \text{if } X[i] \ne Y[j] \text{ (Substitution)}\\
			0 & \text{otherwise}
		\end{cases}
	\end{cases}
	\]
	
\end{frame}


\begin{frame}
	\frametitle{\insertshortsubtitle: \insertsection}
	\framesubtitle{\insertsubsection: Levenshtein distance (Example)}

	\begin{exampleblock}{Example of Levenshtein distance calculation}
		\centering\scriptsize
		\begin{tblr}{
				colspec = {|Q[c,m,wd=3em]|Q[c,m,wd=3em]|c|c|c|c|c|c|c|c|},
				row{1} = {bg=darkblue, fg=white, font=\bfseries},
				column{1} = {bg=darkblue, fg=white, font=\bfseries},
				stretch=1.5,
				colsep=2pt,
			} 	
			
			\hline
			& \# & i &\bfseries m &\bfseries m & a & t & u & r & e \\
			\hline
			\bfseries \# & \SetCell{bg=green!70} 0 & $ \leftarrow $ 1 & $ \leftarrow $ 2 & $ \leftarrow $ 3 & $ \leftarrow $ 4 & $ \leftarrow $ 5 & $ \leftarrow $ 6 & $ \leftarrow $ 7 & $ \leftarrow $ 8\\
			\hline
			\bfseries a & \SetCell{bg=green!70} $ \uparrow $ 1 & $ \nwarrow\leftarrow\uparrow $ 2 & $ \nwarrow\leftarrow\uparrow $ 3 & $ \nwarrow\leftarrow\uparrow $ 4 & $ \nwarrow $ 3 & $ \leftarrow $ 4 & $ \leftarrow $ 5 & $ \leftarrow $ 6 & $ \leftarrow $ 7 \\
			\hline
			\bfseries m & \SetCell{bg=green!70} $ \uparrow $ 2 & $ \nwarrow\leftarrow\uparrow $ 3 & $\nwarrow $ 2 & $\nwarrow\leftarrow $ 3 & $\leftarrow\uparrow $ 4 & $\nwarrow\leftarrow\uparrow $ 5 & $\nwarrow\leftarrow\uparrow $ 6 & $\nwarrow\leftarrow\uparrow $ 7 & $\nwarrow\leftarrow\uparrow $ 8\\
			\hline
			\bfseries i & $ \uparrow $ 3 & \SetCell{bg=green!70} $ \nwarrow $ 2 & \SetCell{bg=green!70} $\leftarrow\uparrow $ 3 & \SetCell{bg=green!70} $\nwarrow\leftarrow\uparrow $ 4 & \SetCell{bg=green!70} $\nwarrow\leftarrow\uparrow $ 5 & \SetCell{bg=green!70} $\nwarrow\leftarrow\uparrow $ 6 & \SetCell{bg=green!70} $\nwarrow\leftarrow\uparrow $ 7 & $\nwarrow\leftarrow\uparrow $ 8 & $\nwarrow\leftarrow\uparrow $ 9\\
			\hline
			\bfseries n & $ \uparrow $ 4 & $ \uparrow $ 3 & $\nwarrow\leftarrow\uparrow $ 4 & $\nwarrow\leftarrow\uparrow $ 5 & $\nwarrow\leftarrow\uparrow $ 6 & $\nwarrow\leftarrow\uparrow $ 7 & $\nwarrow\leftarrow\uparrow $ 8 & \SetCell{bg=green!70} $\nwarrow\leftarrow\uparrow $ 9 & $\nwarrow\leftarrow\uparrow $ 10\\
			\hline
			\bfseries e & $ \uparrow $ 5 & $ \uparrow $ 4 & $\nwarrow\leftarrow\uparrow $ 5 & $\nwarrow\leftarrow\uparrow $ 6 & $\nwarrow\leftarrow\uparrow $ 7 & $\nwarrow\leftarrow\uparrow $ 8 & $\nwarrow\leftarrow\uparrow $ 9 & $\nwarrow\leftarrow\uparrow $ 10 & \SetCell{bg=green!70} $\nwarrow $ 9\\
			\hline
		\end{tblr}
	\end{exampleblock}

\end{frame}

\begin{frame}
	\frametitle{\insertshortsubtitle: \insertsection}
	\framesubtitle{\insertsubsection: Some applications}
	
	\begin{itemize}
		\item \optword{File revision}: e.g., the UNIX \expword{diff} command, which compares two files.
		\item \optword{Spell checking}: suggesting possible corrections for spelling errors (e.g., \expword{Hunspell}).
		\item \optword{Plagiarism detection}: edit distance can be applied at the word level instead of the character level.
		\item \optword{Spam filtering}: spammers may intentionally introduce spelling errors to evade spam detection tools.
		\item \optword{Bioinformatics}: measuring similarity between DNA sequences.
	\end{itemize}
	
\end{frame}

\begin{frame}
	\frametitle{\insertshortsubtitle: \insertsection}
	\framesubtitle{\insertsubsection: Some humor}

	\begin{center}
		\vgraphpage[.4\textheight]{humor/humor-spell.jpg}
		\vgraphpage[.4\textheight]{humor/humor-spell1.jpg}
		\vgraphpage[.4\textheight]{humor/humor-spell2.jpg}
	\end{center}

\end{frame}

%===================================================================================
\section{Text segmentation}
%===================================================================================

\begin{frame}
	\frametitle{\insertshortsubtitle}
	\framesubtitle{\insertsection}
	
	\begin{itemize}
		\item Paragraph boundaries are typically indicated by line breaks or HTML tags such as \keyword{\textless p\textgreater}.
		\begin{itemize}
			\item Text extracted from PDFs often contains extra line breaks. 
		\end{itemize}
		\item Sentence boundaries are usually indicated by periods or other punctuation marks.
		\begin{itemize}
			\item Sentence-ending punctuation is not always reliable (e.g., abbreviations). 
		\end{itemize}
		\item Words are generally separated by spaces.
		\begin{itemize}
			\item Word boundaries are not always indicated by spaces (e.g., in Chinese, Japanese, or URLs).
		\end{itemize}
	\end{itemize}
	
\end{frame}


\subsection{Sentence boundary disambiguation}

\begin{frame}
	\frametitle{\insertshortsubtitle: \insertsection}
	\framesubtitle{\insertsubsection: Problem} 
	
	\begin{itemize}
		\item \expword{/[.?!]/} is a simple regular expression to delimit sentences in languages such as French and English.
		\item Periods are also used in numbers; e.g., \expword{123,456.78} (American style) and \expword{123.456,78} (European style).
		\item Abbreviations contain periods. Examples: \expword{exempli gratia (e.g.), Doctor of Philosophy (Ph.D.), page (p.)}.
		\item Long sentences are difficult to process; it is often useful to split them. Example: \expword{After he woke up, Karim showered, brewed a cup of coffee, and settled down to start his workday.}.
		\item Quotations can complicate sentence splitting. Example: \expword{After receiving the news, Karim exclaimed, "I can't believe it! This is incredible!"}
		\item Some languages, e.g., Thai, do not use explicit markers to delimit sentences.
	\end{itemize}

\end{frame}

\begin{frame}
	\frametitle{\insertshortsubtitle: \insertsection}
	\framesubtitle{\insertsubsection: Features} 

	\begin{itemize}
		\item \optword{Letter case}: sentences typically start with an uppercase letter.
		\item \optword{Parts of speech (POS)}: POS tags can help determine sentence boundaries.
		\item \optword{Word length}: abbreviations are usually short.
		\item \optword{Prefixes and suffixes}: words with affixes are less likely to be abbreviations.
		\item \optword{Abbreviation's class}: certain abbreviations can occur at the end of a sentence (e.g., \expword{Inc.}).
		\item \optword{Proper nouns}: proper nouns start with uppercase letters but may appear mid-sentence, potentially causing ambiguity.
	\end{itemize}

\end{frame}

\begin{frame}
	\frametitle{\insertshortsubtitle: \insertsection}
	\framesubtitle{\insertsubsection: Detection (Rule based)} 
	
	\vskip-3pt
	\begin{exampleblock}{Example of a simple rule-based SBD algorithm}
		\tiny
		\begin{algorithm}[H]
			\vskip-3pt
			\KwData{text}
			\KwResult{list of sentences' start positions}
			
			positions = [0]\;
			
			\ForEach{delimiter $ d_{\sim !?.} \in text $ }{
				
				\eIf{$ d_{\sim !?} $}{
					positions \textleftarrow\ positions $ \cup $ \{position(d)+1\}
				}{
					\If{word after d starts with uppercase}{
						\eIf{word before d has infixes or has more than 5 letters or not a noun}{
							positions \textleftarrow\ positions $ \cup $ \{position(d)+1\}
						}{
							\eIf{word before d not in list of abbreviations}{
								positions \textleftarrow\ positions $ \cup $ \{position(d)+1\}
							}{
								\If{word before d $ \in $ list of ending abbr.}{positions \textleftarrow\ positions $ \cup $ \{position(d)+1\}}
							}
						}
					}
				}
			}
			
			\Return positions\;
			
		\end{algorithm}
	\end{exampleblock}
	
\end{frame}

\begin{frame}
	\frametitle{\insertshortsubtitle: \insertsection}
	\framesubtitle{\insertsubsection: Detection (ML-based)} 
	
	\begin{figure}
		\centering
		\hgraphpage[0.9\textwidth]{SBD-basicML.pdf}
		\caption{Example of ML-based SBD model with predifined features}
	\end{figure}
	
\end{frame}

\begin{frame}
	\frametitle{\insertshortsubtitle: \insertsection}
	\framesubtitle{\insertsubsection: Some humor} 


	\begin{center}
		\vgraphpage{humor/humor-segmentation.jpg}
	\end{center}

\end{frame}

\subsection{Word tokenization}

\begin{frame}
	\frametitle{\insertshortsubtitle: \insertsection}
	\framesubtitle{\insertsubsection: Problem} 

	\begin{itemize}
		\item \expword{/[ ]+/} is a regular expression commonly used to split words (Arabic, French, English, etc.).
		\item Some languages, such as Japanese, do not use explicit word boundaries (\expword{\begin{CJK}{UTF8}{min}今年は本当に忙しかったです。\end{CJK}}).
		\item Compound words can be:
		\begin{itemize}
			\item Open: ``\expword{wake up}''  
			\item Closed: ``\expword{makeup}'', German: ``\expword{Lebensversicherung}'' (life insurance), Arabic: ``\expword{\RL{y_htbrwnhm}}'' (they test them)  
			\item Hyphenated: ``\expword{well-prepared}'', French: ``\expword{va-t-il}'' (will he), ``\expword{c-à-dire}'' (which means)
		\end{itemize}
		\item Ambiguity with certain characters: apostrophes may indicate quotations or contractions (e.g., \expword{She's}, \expword{J'ai}).
		\item Multi-word expressions: e.g., numerical expressions such as dates.
	\end{itemize}

\end{frame}

\begin{frame}
	\frametitle{\insertshortsubtitle: \insertsection}
	\framesubtitle{\insertsubsection: Approaches}
	
	\begin{itemize}
		\item \optword{Rule-based approach}: employing regular expressions (REs) to identify word boundaries.
		\begin{itemize}
			\item Examples and tools:
			\begin{itemize}
				\item \url{https://www.nltk.org/api/nltk.tokenize.html}
				\item \url{https://nlp.stanford.edu/software/tokenizer.shtml}
				\item \url{https://spacy.io/}
				\item \url{https://github.com/kariminf/jslingua}
				\item \url{https://github.com/linuxscout/pyarabic}
			\end{itemize}
		\end{itemize}
		\item \optword{Statistical approach}: using probabilistic or language models to estimate the likelihood that a character indicates a word boundary.
		\begin{itemize}
			\item Examples and tools:
			\begin{itemize}
				\item \url{https://nlp.stanford.edu/software/segmenter.html}
				\item \url{https://opennlp.apache.org/}
			\end{itemize}
		\end{itemize}
	\end{itemize}

\end{frame}

\begin{frame}
	\frametitle{\insertshortsubtitle: \insertsection}
	\framesubtitle{\insertsubsection: Some humor}

	\begin{center}
		\vgraphpage{humor/humor-tokenization.png}
	\end{center}

\end{frame}


%===================================================================================
\section{Text normalization and filtering}
%===================================================================================

\begin{frame}
	\frametitle{\insertshortsubtitle}
	\framesubtitle{\insertsection}
	
	\begin{itemize}
		\item \optword{Text normalization}
		\begin{itemize}
			\item \textbf{Problem}: Texts may contain multiple variations of the same term. For certain tasks, such as information retrieval, the exact form of the text may be unnecessary.
			\item \textbf{Solution}: Transform the text into a canonical form (e.g., lowercasing, stemming/lemmatization, accent removal).
			\item Example challenge: \url{https://www.kaggle.com/c/text-normalization-challenge-english-language}
		\end{itemize}
		\item \optword{Text filtering}
		\begin{itemize}
			\item \textbf{Problem}: Texts may contain non-informative characters, words, or expressions.
			\item \textbf{Solution}: Remove or filter out these elements (e.g., punctuation, stopwords, emojis).
		\end{itemize}
	\end{itemize}

\end{frame}

\subsection{Text normalization}

\begin{frame}
	\frametitle{\insertshortsubtitle: \insertsection}
	\framesubtitle{\insertsubsection\ (1)}
	
	\begin{itemize}
		\item Acronyms and abbreviations:
		\begin{itemize}
			\item Standard form: \expword{US \textrightarrow\ USA, U.S.A. \textrightarrow\ USA}
			\item Long version: \expword{Mr. \textrightarrow\ Mister}
		\end{itemize}
		
		\item Standardizing formats for values such as dates and numbers:
		\begin{itemize}
			\item Conversion to textual form: \expword{1205 DZD \textrightarrow\ One thousand, two hundred five Algerian dinars}.
			\item Unified date format: \expword{12 janvier 1986, 12.01.86, January 12th, 1986 \textrightarrow\ 1986-01-12}.
			\item Replacement by semantic type: \expword{January 12th, 1986 \textrightarrow\ DATE, kariminfo0@gmail.com \textrightarrow\ EMAIL}.
		\end{itemize}
		
		\item Convert uppercase letters to lowercase:
		\begin{itemize}
			\item \expword{Text \textrightarrow\ text}
			\item Exceptions: proper nouns may retain capitalization (e.g., \expword{Will})
		\end{itemize}
	\end{itemize}

\end{frame}

\begin{frame}
	\frametitle{\insertshortsubtitle: \insertsection}
	\framesubtitle{\insertsubsection\ (2)}
	
	\begin{itemize}
		\item Contractions:
		\begin{itemize}
			\item Expansion of contractions: \expword{y'll \textrightarrow\ you all, s'il \textrightarrow\ si il}
		\end{itemize}
		
		\item Diacritics:
		\begin{itemize}
			\item Removal of accents:  \expword{système \textrightarrow\ systeme}
			\item Removal of vowel diacritics (unvocalization):  \expword{\RL{yadrusu} \textrightarrow\ \RL{ydrs}}  
			\item Diacritics may be retained when required (e.g., poetry)
		\end{itemize}
		
		\item Encoding:
		\begin{itemize}
			\item A consistent text encoding should be used during processing (e.g., UTF-8)
		\end{itemize}
		
	\end{itemize}

\end{frame}

\subsection{Text filtering}

\begin{frame}
	\frametitle{\insertshortsubtitle: \insertsection}
	\framesubtitle{\insertsubsection}
	
	\begin{itemize}
		\item Special characters:
		\begin{itemize}
			\item Non-printable or unusual characters can cause errors during text processing.
		\end{itemize}
		
		\item Markup tags:
		\begin{itemize}
			\item Tags such as HTML, XML, or other markup languages should be removed.
		\end{itemize}
		
		\item \keyword{Stop words}:
		\begin{itemize}
			\item Insignificant words such as prepositions, articles, and pronouns.
			\item In information retrieval, these words are typically not indexed.
			\item In automatic text summarization, stopwords may affect sentence scoring.
		\end{itemize}
	\end{itemize}

\end{frame}

\begin{frame}
	\frametitle{\insertshortsubtitle: \insertsection}
	\framesubtitle{\insertsubsection: Some humor}

	\begin{center}
		\vgraphpage{humor/humor-stopwords.jpg}
	\end{center}

\end{frame}

%===================================================================================
\section{Morphology}
%===================================================================================

\begin{frame}
	\frametitle{\insertshortsubtitle}
	\framesubtitle{\insertsection}

	\begin{itemize}
		\item Some languages form words using:
		\begin{itemize}
			\item Inflection (e.g., \expword{conjugation})
			\item Derivation (e.g., \expword{nominalization})
		\end{itemize}
		
		\item Affixation is the most common morphological process.
		
		\item Automating word formation can support tasks such as natural language generation (NLG).
		
		\item Conversely, morphological analysis (decomposition) can assist tasks such as natural language understanding (NLU).
	\end{itemize}

\end{frame}

\subsection{Word formation}

\begin{frame}
	\frametitle{\insertshortsubtitle: \insertsection}
	\framesubtitle{\insertsubsection}
	
	\begin{itemize}
		\item \optword{Inflection} : morphological variation of a word to encode grammatical features (number, gender, tense, etc.).
		\begin{itemize}
			\item Verb conjugation.
			\item Declension of nouns, pronouns, adjectives, adverbs, and articles. 
			\item E.g., Arabic: \expword{\RL{qi.t.t} \textrightarrow\ \RL{qi.ta.t}}.
		\end{itemize}
		
		\item \optword{Derivation}: morphological modification to create a new lexeme or change lexical category.
		\begin{itemize}
			\item Create a new lexeme: \expword{couper \textrightarrow\ découper, \RL{`ml} \textrightarrow\ \RL{ist`ml}}.
			\item Change category: 
			\begin{itemize}
				\item Nominalization: \expword{classer \textrightarrow\ classement, classeur ; \RL{darasa} \textrightarrow\ \RL{darsuN, madrasaTuN, mudarrisuN, dArisuN}}
				\item Adjective formation: \expword{fatiguer \textrightarrow\ fatigant}
			\end{itemize}
		\end{itemize}
	\end{itemize}

\end{frame}

\begin{frame}
	\frametitle{\insertshortsubtitle: \insertsection}
	\framesubtitle{\insertsubsection: Example (Automatic conjugation)}
	
	\begin{itemize}
		\item \optword{Database-based approach}:
		\begin{itemize}
			\item Store verb conjugations directly in a database.
		\end{itemize}
		
		\item \optword{Template-based approach}:
		\begin{itemize}
			\item Store conjugation patterns of representative verbs as templates and classify other verbs according to these templates.
			\item For example, \expword{French verb conjugation}.
		\end{itemize}
		
		\item \optword{Rule-based approach}:
		\begin{itemize}
			\item Use IF-ELSE rules and regular expressions.
			\item Examples: Arabic, English, French, and Japanese verbs conjugation: \url{https://github.com/kariminf/jslingua}
			\item Arabic verbs conjugation: \url{https://github.com/linuxscout/qutrub}
		\end{itemize}
	\end{itemize}

\end{frame}

\begin{frame}
	\frametitle{\insertshortsubtitle: \insertsection}
	\framesubtitle{\insertsubsection: Some humor}

	\begin{center}
		\vgraphpage{humor/humor-formation.jpg}
	\end{center}

\end{frame}

\subsection{Form reducing}

\begin{frame}
	\frametitle{\insertshortsubtitle: \insertsection}
	\framesubtitle{\insertsubsection: Stemming}
	
	\begin{itemize}
		\item Reduction of word forms by removing affixes.
		\item The result is called the \keyword{stem}.
		\item Examples: \expword{chercher \textrightarrow\ cherch}, \expword{arguing \textrightarrow\ argu}.
		\item \optword{Database-based approach}: store all terms and their corresponding stems in a table.
		\item \optword{Statistical approach}: use an n-gram language model to estimate the appropriate truncation position.
		\item \optword{Rule-based approach}:
		\begin{itemize}
			\item Porter algorithm \cite{1980-porter}.
			\item Arabic stemmer: \url{https://github.com/assem-ch/arabicstemmer}
			\item Snowball algorithms: \url{https://snowballstem.org/algorithms/}
		\end{itemize}
	\end{itemize}

\end{frame}

\begin{frame}
	\frametitle{\insertshortsubtitle: \insertsection}
	\framesubtitle{\insertsubsection: Stemming (Porter algorithm)}
	
	\begin{itemize}
		\item The Porter stemming algorithm consists of a set of rules, each with a condition and an action.
		\item \url{https://snowballstem.org/}: a framework to create stemmers.
		\item \optword{Condition on the stem}:
		\begin{itemize}
			\item Examples: stem length, presence of vowels, specific endings.
			\item E.g., \expword{(*v*) Y \textrightarrow\ I}: \expword{happy \textrightarrow\ happi}, \expword{sky \textrightarrow\ sky} 
			(*v* indicates the stem contains at least one vowel)
		\end{itemize}
		\item \optword{Condition on the affix}:
		\begin{itemize}
			\item Rules are applied only to suffixes.
			\item E.g., \expword{SSES \textrightarrow\ SS, ATIONAL \textrightarrow\ ATE}
		\end{itemize}
		\item \optword{Condition on the rules}:
		\begin{itemize}
			\item If a rule is applied, subsequent rules in the same step are not executed.
		\end{itemize}
	\end{itemize}

\end{frame}

\begin{frame}
	\frametitle{\insertshortsubtitle: \insertsection}
	\framesubtitle{\insertsubsection: Lemmatization}
	
	\begin{itemize}
		\item Find the canonical (dictionary) form of a word.
		\item The result is called the \keyword{lemma}.
		\item Examples:
		\begin{itemize}
			\item \expword{comprennent \textrightarrow\ comprendre}, \expword{better \textrightarrow\ good}
			\item Context-dependent: \expword{saw (V) \textrightarrow\ see}, \expword{saw (N) \textrightarrow\ saw}
		\end{itemize}
		\item Lemmatization uses the word's context and part of speech.
		\item \optword{Lexical resources}:
		\begin{itemize}
			\item \url{https://www.nltk.org/api/nltk.stem.html} (NLTK WordNetLemmatizer)
			\item \url{https://github.com/sloria/textblob}
			\item \url{https://spacy.io/}
		\end{itemize}
		\item \optword{Machine learning approaches}:
		\begin{itemize}
			\item \url{https://opennlp.apache.org/}
		\end{itemize}
	\end{itemize}

\end{frame}

\begin{frame}[fragile]
	\frametitle{\insertshortsubtitle: \insertsection}
	\framesubtitle{\insertsubsection: Lemmatization (E.g. Morphy)}

	\vspace{-6pt}
	\begin{block}{Wordnet's ``morphy'' Lemmatization}
		\footnotesize
		\begin{algorithm}[H]
			\vskip-3pt
			\KwData{word, POS}
			\KwResult{list of possible lemmas}
			
			\If{word $ \in $  exceptions\_list[POS]}{
				\Return search\_in\_dictionary(\{word\} $ \cup $ exceptions\_list[POS]) \tcp*{handle irregular forms}
			}
			
			forms = \{word\}\;
			
			\While{forms $ \ne \emptyset $}{
				forms = delete\_affixes(forms, POS)\tcp*{generate candidate forms by removing affixes}
				
				results = search\_in\_dictionary(\{word\} $ \cup $ forms)\;
				
				\lIf{results $ \ne \emptyset $ }{
					\Return results \tcp*{return matches from dictionary}
				}
			}
			
			\Return $ \emptyset $ \tcp*{no lemma found}
	
		\end{algorithm}
	\end{block}

\end{frame}

\begin{frame}
	\frametitle{\insertshortsubtitle: \insertsection}
	\framesubtitle{\insertsubsection: Some humor}

	\begin{center}
		\vgraphpage{humor/humor-stemming.jpeg}
	\end{center}

\end{frame}

\insertbibliography{NLP03}{*}

\end{document}

